\chapter{Sprint 1 : Gestion des utilisateurs}

\section*{Introduction}
Le premier sprint établit les fondations de la plateforme : authentification sécurisée
via Laravel Sanctum, contrôle d'accès par rôles (Spatie Permission), gestion des
profils utilisateurs, préférences de notification, et tableau de bord administrateur
avec KPI temps réel.

% ============================================================
\section{Sprint Backlog}

\begin{table}[H]
  \centering
  \caption{Sprint Backlog — Sprint 1}
  \label{tab:backlog_s1}
  \begin{tabularx}{\textwidth}{C{1.2cm}L{5cm}L{2cm}C{2cm}C{1.5cm}}
    \toprule
    \thead{ID} & \thead{User Story} & \thead{Acteur} & \thead{Priorité} & \thead{Points} \\
    \midrule
    \rowcolor{lightgray}
    US-01 & S'inscrire et se connecter en sécurité              & Tous  & Must   & 5 \\
    US-02 & Gérer les comptes utilisateurs                      & Admin & Must   & 8 \\
    \rowcolor{lightgray}
    US-03 & Définir rôles et permissions Spatie                 & Admin & Must   & 5 \\
    US-04 & Modifier profil et préférences                      & Tous  & Should & 5 \\
    \rowcolor{lightgray}
    US-05 & Réinitialiser mot de passe par email                & Tous  & Should & 3 \\
    \midrule
    \multicolumn{4}{r}{\textbf{Total :}} & \textbf{26} \\
    \bottomrule
  \end{tabularx}
\end{table}

% ============================================================
\section{Conception}

\subsection{Diagramme de cas d'utilisation}

\begin{figure}[H]
  \centering
  \IfFileExists{figures/sprint1_use_case.png}{%
    \includegraphics[width=0.88\textwidth]{figures/sprint1_use_case.png}%
  }{%
    \fbox{\parbox{0.88\textwidth}{\centering\vspace{1.8cm}
      \textbf{[IMAGE A INSÉRER]}\\[0.4em]
      \texttt{figures/sprint1\_use\_case.png}\\[0.3em]
      \textit{Diagramme de cas d'utilisation — Sprint 1 : Gestion des utilisateurs}
    \vspace{1.8cm}}}%
  }
  \caption{Diagramme de cas d'utilisation — Sprint 1}
  \label{fig:uc_s1}
\end{figure}

Les acteurs impliqués sont \textbf{Admin}, \textbf{Tous} (Member/Client)
et le \textbf{Système Email}. Les cas couvrent : S'inscrire, Se connecter
(Sanctum), Se déconnecter, Réinitialiser le mot de passe (avec envoi email),
Modifier son profil, Changer la langue et le thème, Configurer les notifications,
Gérer les spécialités, Gérer les comptes (Admin), Définir les rôles Spatie (Admin)
et Consulter le dashboard administrateur (Admin).

\subsection{Diagramme de séquence — Connexion Sanctum}

\begin{figure}[H]
  \centering
  \IfFileExists{figures/sprint1_sequence_login.png}{%
    \includegraphics[width=0.90\textwidth]{figures/sprint1_sequence_login.png}%
  }{%
    \fbox{\parbox{0.90\textwidth}{\centering\vspace{1.8cm}
      \textbf{[IMAGE A INSÉRER]}\\[0.4em]
      \texttt{figures/sprint1\_sequence\_login.png}\\[0.3em]
      \textit{Séquence de connexion : React → AuthController → User → Sanctum → MySQL}
    \vspace{1.8cm}}}%
  }
  \caption{Diagramme de séquence — Connexion avec Laravel Sanctum}
  \label{fig:seq_login}
\end{figure}

Le flux nominal se déroule en sept étapes : (1) saisie email/mot de passe,
(2) validation Zod côté React, (3) \texttt{POST /api/auth/login},
(4) recherche utilisateur en base, (5) vérification du hash Bcrypt,
(6) création du token Sanctum, (7) réponse \texttt{\{token, UserResource\}}.
Deux alternatives sont gérées : compte désactivé (422) et mauvais mot de passe (422).

\subsection{Diagramme de classes}

\begin{figure}[H]
  \centering
  \IfFileExists{figures/sprint1_class.png}{%
    \includegraphics[width=0.85\textwidth]{figures/sprint1_class.png}%
  }{%
    \fbox{\parbox{0.85\textwidth}{\centering\vspace{1.8cm}
      \textbf{[IMAGE A INSÉRER]}\\[0.4em]
      \texttt{figures/sprint1\_class.png}\\[0.3em]
      \textit{Classes : User, PersonalAccessToken, Role, Permission, Specialty}
    \vspace{1.8cm}}}%
  }
  \caption{Diagramme de classes — Sprint 1}
  \label{fig:class_s1}
\end{figure}

% ============================================================
\section{Réalisation}

\subsection{Page de connexion}

\begin{figure}[H]
  \centering
  \IfFileExists{figures/sprint1_01_login.png}{%
    \includegraphics[width=0.88\textwidth]{figures/sprint1_01_login.png}%
  }{%
    \fbox{\parbox{0.88\textwidth}{\centering\vspace{2cm}
      \textbf{[IMAGE A INSÉRER]}\\[0.4em]
      \texttt{figures/sprint1\_01\_login.png}\\[0.3em]
      \textit{Page de connexion : deux colonnes, panneau violet animé, formulaire React Hook Form + Zod}
    \vspace{2cm}}}%
  }
  \caption{Page de connexion — TAC-TIC IntelliPM}
  \label{fig:ui_login}
\end{figure}

La page de connexion adopte une mise en page deux colonnes. Le panneau gauche présente
une animation de fond avec des orbes colorées (Framer Motion). Le formulaire droit
utilise \textbf{React Hook Form} pour la gestion d'état et \textbf{Zod} pour la
validation côté client, avant l'appel à l'API.

\subsection{Tableau de bord administrateur — KPI}

\begin{figure}[H]
  \centering
  \IfFileExists{figures/sprint1_02_dashboard_kpi.png}{%
    \includegraphics[width=0.88\textwidth]{figures/sprint1_02_dashboard_kpi.png}%
  }{%
    \fbox{\parbox{0.88\textwidth}{\centering\vspace{2cm}
      \textbf{[IMAGE A INSÉRER]}\\[0.4em]
      \texttt{figures/sprint1\_02\_dashboard\_kpi.png}\\[0.3em]
      \textit{5 cartes KPI : 13 utilisateurs, 12 projets, 91 tâches, 41\,\% complétion, 37 en retard}
    \vspace{2cm}}}%
  }
  \caption{Dashboard administrateur — Cartes KPI}
  \label{fig:ui_dashboard_kpi}
\end{figure}

\subsection{Tableau de bord — Charts analytiques}

\begin{figure}[H]
  \centering
  \IfFileExists{figures/sprint1_03_dashboard_charts.png}{%
    \includegraphics[width=0.88\textwidth]{figures/sprint1_03_dashboard_charts.png}%
  }{%
    \fbox{\parbox{0.88\textwidth}{\centering\vspace{2cm}
      \textbf{[IMAGE A INSÉRER]}\\[0.4em]
      \texttt{figures/sprint1\_03\_dashboard\_charts.png}\\[0.3em]
      \textit{Donut répartition statuts, graphe en barres charge équipes, avancement par projet}
    \vspace{2cm}}}%
  }
  \caption{Dashboard administrateur — Charts analytiques}
  \label{fig:ui_dashboard_charts}
\end{figure}

\subsection{Page Paramètres}

\begin{figure}[H]
  \centering
  \begin{subfigure}[b]{0.48\textwidth}
    \IfFileExists{figures/sprint1_04_settings_profil.png}{%
      \includegraphics[width=\textwidth]{figures/sprint1_04_settings_profil.png}%
    }{%
      \fbox{\parbox{\textwidth}{\centering\vspace{1.2cm}\small
        \texttt{sprint1\_04\_settings\_profil.png}\\Onglet Profil\vspace{1.2cm}}}%
    }
    \caption{Onglet Profil}
  \end{subfigure}
  \hfill
  \begin{subfigure}[b]{0.48\textwidth}
    \IfFileExists{figures/sprint1_07_settings_notifs.png}{%
      \includegraphics[width=\textwidth]{figures/sprint1_07_settings_notifs.png}%
    }{%
      \fbox{\parbox{\textwidth}{\centering\vspace{1.2cm}\small
        \texttt{sprint1\_07\_settings\_notifs.png}\\Préférences notifications\vspace{1.2cm}}}%
    }
    \caption{Préférences de notification}
  \end{subfigure}
  \caption{Page Paramètres — Onglets Profil et Notifications}
  \label{fig:ui_settings}
\end{figure}

\subsection{Code clé — \texttt{AuthController::login()}}

\begin{lstlisting}[style=phpstyle,caption={AuthController::login() — Authentification Sanctum},label={lst:auth_login}]
// app/Http/Controllers/Api/AuthController.php
public function login(LoginRequest $request)
{
    $user = User::where('email', $request->email)->first();

    if (!$user || !Hash::check($request->password, $user->password)) {
        throw ValidationException::withMessages(
            ['email' => __('auth.failed')]
        );
    }
    if ($user->status === 'disabled') {
        throw ValidationException::withMessages(
            ['email' => __('auth.disabled')]
        );
    }

    $token = $user->createToken('intellipm-web')->plainTextToken;

    return response()->json([
        'token' => $token,
        'user'  => UserResource::make($user->load('roles', 'specialties')),
    ]);
}
\end{lstlisting}

% ============================================================
\section{Tests unitaires}

\begin{lstlisting}[style=phpstyle,caption={Tests unitaires — Authentification (PestPHP)},label={lst:tests_s1}]
it('rejects disabled accounts', function () {
    $user = User::factory()->create(['status' => 'disabled']);
    $this->postJson('/api/auth/login', [
        'email'    => $user->email,
        'password' => 'password',
    ])->assertStatus(422)
      ->assertJsonValidationErrors(['email']);
});

it('returns token on valid credentials', function () {
    $user = User::factory()->create();
    $this->postJson('/api/auth/login', [
        'email'    => $user->email,
        'password' => 'password',
    ])->assertOk()
      ->assertJsonStructure(['token', 'user']);
});

it('blocks access to protected routes without token', function () {
    $this->getJson('/api/auth/me')
         ->assertUnauthorized();
});
\end{lstlisting}

\section*{Conclusion}
Le sprint 1 a été livré dans les délais impartis. Les fondations de sécurité
sont opérationnelles : authentification Sanctum, RBAC Spatie à deux niveaux,
protection des routes par middleware. Le tableau de bord administrateur a reçu
un accueil très favorable lors de la démonstration de fin de sprint. Le seul
point d'amélioration noté concerne la configuration initiale de Spatie Permission,
légèrement sous-estimée lors du planning.
